\scanpage{6}
\diagram{\begin{tikzcd}[ampersand replacement=\&,row sep=large,column sep=large]
C\times_S B \arrow[r,"\psi\times\operatorname{id}_B"] \arrow[d,"\operatorname{pr}_1"'] \arrow[dr,phantom,"\text{\textcircled{cart}}"] \& B^*\times_S B^{\deleted{\unsure{*}}} \arrow[d,"\operatorname{pr}'_1"] \\
C \arrow[r,"\psi"'] \& B^*
\end{tikzcd}}
Ceci dit, on aura\\
\deleted{$\mathbf R\operatorname{pr}_{2*}(\mathcal N)$}
\[
\mathcal N\simeq(\psi\times\operatorname{id}_B)^*(W_B)
\]
($W_B$ le Module de Weil sur $B\times_S B^*$), d'où
\[
\begin{aligned}
\mathbf R\operatorname{pr}_{2*}(\mathcal N)&\simeq\mathbf R\operatorname{pr}_{2*}(\psi\times\operatorname{id}_B)^*(W_B)\\
&\overset{\text{chgt de base}}{\simeq}\mathbf L\psi^*(\mathbf R\operatorname{pr}'_{2*}(W_B)).
\end{aligned}
\]
\uline{Lemme 1}. Soit $B$ un schéma abélien sur $S$, \inserted{de dim relative $n$,} $W$ le\\
faisceau de Weil sur $B^{\unsure{*}}\times_S B^{\deleted{\unsure{*}}}$, $p:B^*\times_S B^{\deleted{\unsure{*}}}\mathbin{-}B^*$ la\\
projection. Alors
\[
\mathbf Rp^{\deleted{\unsure{*}}}_*(W)\simeq
 e_{\unsure{0|B^*}*}\bigl[\deleted{\unsure{\mathcal O_{B^*}[-n]|\mathcal O_S[-n]}\unsure{\otimes\tau_{B^*}}}\tau_{B^*}[-n]\bigr)
\]
(isom. canonique), où $\tau_{B^*}=\overset{n}{\Lambda}\underline t_{B^*}$ ($t_{B^*}$ le faisceau tangent à\\
$B^*$ le long de la section $0$).
\margin{\figure{\begin{tikzpicture}[x=0.8cm,y=0.8cm]
\draw (0,2.3)..controls (1.4,2.5) and (1.7,2.2)..(1.7,2.0)..controls (1.7,1.6) and (2.5,1.7)..(2.4,1.2)..controls (2.3,0.5) and (0.8,0.6)..(0.2,0);
\end{tikzpicture}}
$\simeq e_*(\omega^{-1})$\\
où $\omega=\unsure{\det^*|\det e^*}(\Omega^1_{\unsure{B^*/S|B^*/A}})$}
On en conclut un

\uuline{Corollaire} Soient $C,B$ schémas abéliens sur $S$, \inserted{$B$ de dim relative $n$,} $\mathcal N$\\
un Module inversible birigidifié sur $C\times_S B$, $\psi:C\longrightarrow B^*$\\
l'hom correspondant. Alors, on a
\[
\mathbf R\operatorname{pr}_{1*}(\mathcal N)\simeq\mathbf L\psi^*\bigl(\tau_{B^*}\deleted{\unsure{*}}[-n]\bigr)
\]
\margin{\deleted{\unsure{a}}}
En particulier, si $\psi$ est un épimorphisme, \deleted{\unsure{alors}}\\
plat, et si $K=\operatorname{Ker}\psi$, alors $K$ est un sous-schéma\\
en groupe de $C$ plat de prés. finie sur $S$, donc \deleted{\unsure{dont}}\\
$\uline{\mathcal O}_K$ est \deleted{\unsure{un module plat}} parfait sur $C$, et on a
\[
\mathbf R\operatorname{pr}_{1*}(\mathcal N)\simeq\uline{\mathcal O}_K[-n]\otimes_S\tau_{B^*}.
\]
Dans le cas qui nous occupe, $C=A\times A$, \deleted{$C^*\simeq A$}\\
$\varphi:A\simeq C^*$, et \deleted{\unsure{donc}} par suite $K$ est l'\uline{anti-diagonale}\\
de $A\times_S A$, noyau de $\pi_A:A\times_S A\longrightarrow A$, qui est le\\
graphe de la \uline{symétrie} $s_A$ de $A$. Donc \deleted{\unsure{le même résultat}}\\
modulo dém. du Lemme 1.
\mspage{5}
